% AA-41 | Audyt I: sprawdzenie matematyczne, językowe i przykład w dodatniej charakterystyce.
\chapter{Reprezentacje grup skończonych i twierdzenie Maschkego}
\markboth{Reprezentacje i twierdzenie Maschkego}{Reprezentacje i twierdzenie Maschkego}
\label{aa:ch:representations}
\section{Od działań grup do operatorów liniowych}
Działanie grupy na zbiorze zachowuje jedynie strukturę zbiorową.
Gdy elementy zbioru tworzą przestrzeń wektorową, możemy wymagać,
aby każdy element grupy działał liniowo. Wówczas pytania o niezmiennicze
podprzestrzenie stają się pytaniami o rozkład reprezentacji.
\begin{definicja}[Reprezentacja i podreprezentacja]\label{aa:def:representation}
Niech $G$ będzie grupą, $K$ ciałem, a $V$ przestrzenią wektorową nad $K$.
Reprezentacją $G$ na $V$ nazywamy homomorfizm grup
$\rho:G\to\operatorname{GL}_K(V)$, gdzie $\operatorname{GL}_K(V)$
oznacza grupę wszystkich odwracalnych $K$-liniowych operatorów na $V$.
Podprzestrzeń $W\subseteq V$ jest $G$-niezmiennicza, jeżeli
$\rho(g)w\in W$ dla wszystkich $g\in G$ i $w\in W$.
Reprezentacja jest nieprzywiedlna, gdy $V\ne0$ i jedynymi
podprzestrzeniami niezmienniczymi są $0$ oraz $V$.
\end{definicja}
Warunek homomorfizmu oznacza jawnie $\rho(gh)=\rho(g)\rho(h)$
oraz $\rho(e)=\operatorname{id}_V$. Ponieważ $g^{-1}\in G$,
$\rho(g)^{-1}=\rho(g^{-1})$; zatem $\rho(g)W\subseteq W$
implikuje również $\rho(g)W=W$.
\begin{przyklad}[Reprezentacja znaku]\label{aa:prz:sign-rep}
Niech $G=S_3$, $K=\mathbb Q$, $V=\mathbb Q$.
Definiujemy $\rho(\sigma)(v)=\operatorname{sgn}(\sigma)v$.
Ponieważ znak iloczynu permutacji jest iloczynem znaków,
$\rho$ jest homomorfizmem. Każda podprzestrzeń jednowymiarowego
$V$ jest zerowa albo równa $V$, zatem reprezentacja jest
nieprzywiedlna. Nie jest trywialna: transpozycja działa jako $-1$.
\end{przyklad}
\begin{przyklad}[Reprezentacja permutacyjna i rozkład]\label{aa:prz:permutation-rep}
Niech $S_3$ permutuje bazę $e_1,e_2,e_3$ przestrzeni $\mathbb Q^3$.
Wówczas $L=\mathbb Q(e_1+e_2+e_3)$ jest niezmiennicza.
Także $W=\{(x_1,x_2,x_3):x_1+x_2+x_3=0\}$ jest niezmiennicza,
gdyż permutacja nie zmienia sumy współrzędnych.
Jest ona dwuwymiarowa: wektory $a=e_1-e_2$ i
$b=e_2-e_3$ leżą w $W$ i tworzą jej bazę, ponieważ
$(u,v,w)=(u)a+(-w)b$ dla każdego $(u,v,w)\in W$.
Każdy $v=(x_1,x_2,x_3)$ zapisujemy jako
$v=\frac{x_1+x_2+x_3}{3}(1,1,1)+w$, gdzie $w\in W$.
Ponadto $L\cap W=0$: jeżeli $t(1,1,1)\in W$, to
$3t=0$ w $\mathbb Q$, a więc $t=0$.
Stąd $\mathbb Q^3=L\oplus W$.
Pokażemy też, że podreprezentacja na $W$ jest
nieprzywiedlna nad $\mathbb Q$. Trójcykl
$g=(1\;2\;3)$ przekształca $a$ w $b$, a $b$ w
$-a-b$. Macierz $\rho(g)|_W$ w bazie $(a,b)$ ma więc
postać
\[
\begin{pmatrix}0&-1\\1&-1\end{pmatrix},
\qquad \det(TI-\rho(g)|_W)=T^2+T+1.
\]
Ten wielomian nie ma pierwiastka wymiernego: jego
wyróżnik jest równy $-3$, a kwadrat liczby wymiernej
nie może być ujemny. Gdyby $W$ zawierała niezerową
właściwą podprzestrzeń $S_3$-niezmienniczą, byłaby to
prosta wymierna niezmiennicza również względem $g$.
Operator na tej prostej działałby przez skalar z
$\mathbb Q$, będący wartością własną powyższej macierzy;
takiego skalara nie ma. Zatem $W$ jest nieprzywiedlna.
Widzimy konkretnie sumę reprezentacji wymiarów $1$ i $2$,
a nie tylko dwa abstrakcyjne składniki.
\end{przyklad}
\section{Uśrednianie projekcji}
Jeżeli $W$ jest niezmiennicza, dowolne liniowe dopełnienie
$U$ nie musi być niezmiennicze. Dla grupy skończonej
można poprawić rzut na $W$, uśredniając jego sprzężenia.
\begin{twierdzenie}[Maschke]\label{aa:tw:maschke}
Niech $G$ będzie grupą skończoną, $K$ ciałem takim, że
$\operatorname{char}K$ nie dzieli $|G|$, a $\rho:G\to
\operatorname{GL}_K(V)$ reprezentacją na skończenie wymiarowej
przestrzeni $V$. Dla każdej podprzestrzeni $G$-niezmienniczej
$W\subseteq V$ istnieje podprzestrzeń $G$-niezmiennicza $U$
taka, że $V=W\oplus U$. W konsekwencji każda taka reprezentacja
jest sumą prostą reprezentacji nieprzywiedlnych.
\end{twierdzenie}
\begin{proof}
Wybierzmy bazę $W$ i uzupełnijmy ją do bazy $V$.
Rzut liniowy $p:V\to W$ określony przez $p(w)=w$ dla $w\in W$
oraz $p(u)=0$ na pozostałych wektorach bazy istnieje i spełnia
$p|_W=\operatorname{id}_W$.
Ponieważ $|G|\cdot1_K\ne0$, element $|G|\cdot1_K$ ma
odwrotność w $K$. Określmy
\[
P=\frac1{|G|}\sum_{g\in G}\rho(g)p\rho(g)^{-1}.
\]
Każdy składnik ma obraz w $W$, ponieważ $p(V)\subseteq W$
i $\rho(g)W=W$. Stąd $P(V)\subseteq W$. Ponadto suma jest
$K$-liniowa, ponieważ każdy operator
$\rho(g)p\rho(g)^{-1}$ jest $K$-liniowy.
Jeśli $w\in W$, to $\rho(g)^{-1}w\in W$, więc
$\rho(g)p\rho(g)^{-1}w=w$. Sumując $|G|$ równych składników,
dostajemy $P(w)=w$. W szczególności $P^2=P$.
Dla $h\in G$ przekształcamy sumę przez podstawienie $k=hg$:
\[
\rho(h)P\rho(h)^{-1}=
\frac1{|G|}\sum_{g\in G}\rho(hg)p\rho(hg)^{-1}=P.
\]
Zatem $P\rho(h)=\rho(h)P$. Jądro $U=\ker P$ jest niezmiennicze:
$P(\rho(h)u)=\rho(h)P(u)=0$ dla $u\in U$.
Dla każdego $v\in V$ mamy $v=P(v)+(v-P(v))$, przy czym
$P(v)\in W$ oraz $P(v-P(v))=P(v)-P^2(v)=0$.
Jeżeli $w\in W\cap U$, to $w=P(w)=0$.
To dowodzi $V=W\oplus U$.

Aby uzyskać ostatnią tezę, prowadzimy indukcję względem
$\dim_KV$. Wymiar $0$ daje pustą sumę. Jeżeli $V\ne0$
i jest nieprzywiedlna, stanowi żądaną sumę. W przeciwnym
razie wybieramy niezerową właściwą podprzestrzeń niezmienniczą
$W$. Jej dopełnienie $U$ jest niezmiennicze na mocy pierwszej
części dowodu. Obie przestrzenie mają wymiar mniejszy od
$\dim V$, więc stosujemy do nich założenie indukcyjne.
\end{proof}
\begin{przyklad}[Dlaczego założenie o charakterystyce jest konieczne]
\label{aa:prz:maschke-counter}
Niech $K=\mathbb F_2$, $G=C_2=\{e,s\}$ i $V=K^2$.
Definiujemy $\rho(s)=\begin{pmatrix}1&1\\0&1\end{pmatrix}$.
W charakterystyce $2$ kwadrat tej macierzy jest jednostkowy,
więc otrzymujemy reprezentację $C_2$. Prosta $W=Ke_1$
jest niezmiennicza. Gdyby istniało niezmiennicze dopełnienie,
musiałoby być prostą $K(e_2+ae_1)$ dla pewnego $a\in K$.
Jednak $\rho(s)(e_2+ae_1)=e_2+(a+1)e_1$ nie należy do tej
samej prostej: współczynnik przy $e_2$ wymusza skalar $1$,
a wtedy $a+1\ne a$. Dopełnienie nie istnieje.
\end{przyklad}
\begin{przyklad}[Twierdzenie Maschkego w dodatniej charakterystyce]
\label{aa:prz:maschke-f2-c3}
Niech $K=\mathbb F_2$, $G=C_3=\langle g\rangle$ i niech
$g$ działa na $V=K^3$ przez $g e_1=e_2$, $g e_2=e_3$,
$g e_3=e_1$. Ponieważ $|G|=3$, liczba $3\cdot1_K=1_K$
jest odwracalna, mimo że charakterystyka $K$ jest dodatnia.
Podprzestrzenie
\[
L=K(1,1,1),\qquad
W=\{(a,b,c)\in K^3:a+b+c=0\}
\]
są $G$-niezmiennicze. Mają wymiary odpowiednio $1$ i $2$.
Jeżeli $t(1,1,1)\in W$, to $3t=t=0$ w $K$,
więc $L\cap W=0$ i $V=L\oplus W$.
Zdefiniujmy średnią operatorów działania
$R_G=\frac1{|G|}\sum_{h\in G}\rho(h)$. Ponieważ
$1/3=1$ w $\mathbb F_2$, mamy
\[
R_G(a,b,c)=(a,b,c)+(c,a,b)+(b,c,a)
=(a+b+c)(1,1,1).
\]
Jest to rzut na $L$ wzdłuż $W$.

Pokażmy ponadto, że $W$ jest nieprzywiedlna. Wektory
$a=e_1+e_2$ i $b=e_2+e_3$ tworzą bazę $W$.
Ponieważ $g a=b$ i $g b=e_3+e_1=a+b$, macierz
$g|_W$ w bazie $(a,b)$ wynosi
\[
\begin{pmatrix}0&1\\1&1\end{pmatrix},
\qquad \det(TI-g|_W)=T^2+T+1.
\]
Wielomian $T^2+T+1$ nie ma pierwiastka w $\mathbb F_2$,
bo przy $T=0$ i $T=1$ przyjmuje wartość $1$.
Gdyby $W$ zawierała niezerową właściwą podprzestrzeń
$G$-niezmienniczą, byłaby nią prosta stabilna względem $g$,
a więc istniałaby wartość własna w $K$; sprzeczność.
Otrzymaliśmy zatem rozkład na reprezentacje nieprzywiedlne
w charakterystyce $2$, która nie dzieli rzędu grupy $C_3$.
\end{przyklad}

\begin{stwierdzenie}[Dwie projekcje związane z inwolucją]
\label{aa:prop:involution-projectors}
Niech $K$ będzie ciałem, w którym $2\cdot1_K\ne0$, niech
$V$ będzie skończenie wymiarową przestrzenią wektorową nad $K$,
a $T\in\operatorname{GL}_K(V)$ operatorem spełniającym $T^2=I$.
Wtedy operatory $P_+=(I+T)/2$ i $P_-=(I-T)/2$ są rzutami,
a przestrzeń $V$ jest sumą prostą przestrzeni
$V_+=\ker(T-I)$ oraz $V_-=\ker(T+I)$.
Dokładniej, $\operatorname{im}P_+=V_+$,
$\operatorname{im}P_-=V_-$ i dla każdego $v\in V$ zachodzi
$v=P_+v+P_-v$.
\end{stwierdzenie}
\begin{proof}
Rachunki przeprowadzamy w algebrze operatorów liniowych
na $V$. Z równości $T^2=I$ otrzymujemy
$(I\pm T)^2=I\pm2T+T^2=2(I\pm T)$.
Ponieważ $2$ ma odwrotność w $K$, wynika stąd
$P_+^2=P_+$ oraz $P_-^2=P_-$.
Ponadto $P_++P_-=I$ i
$P_+P_-=(I-T^2)/4=0$.
Dla każdego $v$ mamy $T(P_+v)=P_+v$ i
$T(P_-v)=-P_-v$, gdyż $T$ komutuje z wielomianami od $T$.
Stąd $\operatorname{im}P_+\subseteq V_+$ i
$\operatorname{im}P_-\subseteq V_-$.
Jeżeli $w\in V_+$, to $P_+w=(w+w)/2=w$;
analogicznie $P_-w=w$ dla $w\in V_-$.
Otrzymujemy równości obrazów.
Z $P_++P_-=I$ wynika $V=V_++V_-$.
Gdyby $w\in V_+\cap V_-$, to zarazem $Tw=w$ i
$Tw=-w$, a zatem $2w=0$ i $w=0$.
Suma jest więc prosta.
\end{proof}
\begin{przyklad}[Jawna dekompozycja działania $C_2$]
\label{aa:prz:involution-projectors}
Dla macierzy $T=\begin{pmatrix}0&1\\1&0\end{pmatrix}$
na $\mathbb Q^2$ otrzymujemy
$P_+=\tfrac12\begin{pmatrix}1&1\\1&1\end{pmatrix}$ i
$P_-=\tfrac12\begin{pmatrix}1&-1\\-1&1\end{pmatrix}$.
Pierwszy operator wysyła $(x,y)$ na
$((x+y)/2,(x+y)/2)$, a drugi na
$((x-y)/2,(y-x)/2)$.
Ich obrazy to odpowiednio $\mathbb Q(1,1)$ i
$\mathbb Q(1,-1)$; otrzymujemy w sposób wykonawczy
rozkład z Przykładu~\ref{aa:prz:maschke-explicit}.
W charakterystyce dwa operatory $P_\pm$ nie są
zdefiniowane: nie istnieje element $1/2$ w ciele.
\end{przyklad}

\section{Algebra grupowa}
\begin{definicja}[Algebra grupowa]\label{aa:def:group-algebra-rep}
Dla grupy skończonej $G$ i ciała $K$ algebra $K[G]$ jest
$K$-przestrzenią wektorową z bazą formalnych symboli $[g]$,
$g\in G$, z mnożeniem określonym przez $[g][h]=[gh]$
i przedłużonym dwuliniowo. Elementem jednostkowym jest $[e]$.
\end{definicja}
\begin{stwierdzenie}[Reprezentacje jako moduły]\label{aa:prop:rep-modules}
Dla ustalonych $G,K$ struktury reprezentacji $G$ na
$K$-przestrzeni $V$ są we wzajemnie jednoznacznej odpowiedniości
ze strukturami lewego $K[G]$-modułu na $V$, których działanie
$K$ jest zwykłym mnożeniem przez skalary.
\end{stwierdzenie}
\begin{proof}
Dla reprezentacji $\rho$ określamy
$(\sum_g a_g[g])v=\sum_g a_g\rho(g)v$.
Jest to dobrze określone, ponieważ każdy element $K[G]$
ma jednoznaczne współczynniki. Rozdzielność wynika z liniowości,
a łączność z równości $\rho(g)\rho(h)=\rho(gh)$;
$[e]$ działa identycznościowo. Odwrotnie, z działania
$K[G]$ bierzemy $\rho(g)v=[g]v$. Ponieważ $[g][g^{-1}]=[e]$,
operator $\rho(g)$ jest odwracalny, a reguła mnożenia daje
homomorfizm grup. Obie konstrukcje są wzajemnie odwrotne.
\end{proof}


\begin{przyklad}[Co dzieje się z algebrą grupową w charakterystyce grupy?]
\label{aa:prz:group-algebra-modular}
Niech $G=C_2=\{e,s\}$ i $K=\mathbb F_2$.
W algebrze grupowej $K[G]$ rozpatrzmy
$z=[s]-[e]$. Mamy
\[
z^2=[s]^2-2[s][e]+[e]^2=[e]+[e]=0,
\]
gdyż $s^2=e$ i $2=0$ w $\mathbb F_2$.
Jednocześnie $z\ne0$, ponieważ $[s]$ i $[e]$
są odrębnymi wektorami bazy $K[G]$.
Niezerowy element nilpotentny nie świadczy sam w sobie
o braku półprostoty dowolnego pierścienia nieprzemiennego.
W naszym konkretnym przypadku możemy jednak wykazać
brak niezmienniczego dopełnienia bez takiego wnioskowania.
Lewy moduł regularny $K[G]$ jest dwuwymiarowy, a mnożenie
przez $[s]$ zamienia wektory bazy $[e]$ i $[s]$.
Prosta $W=K([e]+[s])$ jest niezmiennicza, ponieważ
$[s]([e]+[s])=[s]+[e]$.
Każda inna prosta w dwuwymiarowej przestrzeni nad
$\mathbb F_2$ jest równa $K[e]$ albo $K[s]$:
w istocie przestrzeń ma dokładnie trzy wektory niezerowe
$[e]$, $[s]$ i $[e]+[s]$. Żadna z tych dwóch pozostałych
prostych nie jest niezmiennicza względem zamiany bazy.
Stąd $W$ nie ma dopełnienia $G$-niezmienniczego.
Ten rachunek pokazuje mechanizm powstania trudności
w charakterystyce dzielącej rząd grupy i daje
wersję kontrprzykładu~\ref{aa:prz:maschke-counter}
bez użycia macierzy Jordana.
\end{przyklad}

\begin{twierdzenie}[Ostrość warunku twierdzenia Maschkego]
\label{aa:tw:maschke-necessity}
Niech $G$ będzie grupą skończoną i $K$ dowolnym ciałem.
Następujące warunki są równoważne:
\begin{enumerate}
\item $|G|\cdot1_K\ne0$;
\item dla każdej skończenie wymiarowej reprezentacji $G$
nad $K$ i każdej jej podprzestrzeni niezmienniczej
istnieje niezmiennicze dopełnienie;
\item każda skończenie wymiarowa reprezentacja $G$
nad $K$ jest sumą prostą reprezentacji nieprzywiedlnych.
\end{enumerate}
\end{twierdzenie}
\begin{proof}
Implikację $1\Rightarrow2$ daje
Twierdzenie~\ref{aa:tw:maschke}, a $2\Rightarrow3$
jego dowód indukcyjny.
Dla uzasadnienia $3\Rightarrow2$ zauważmy, że każdy
iloraz skończenie wymiarowej reprezentacji, będącej
sumą prostą reprezentacji nieprzywiedlnych, jest ponownie
sumą takich reprezentacji. Istotnie, jeśli
$V=S_1\oplus\cdots\oplus S_r$ z $S_i$ nieprzywiedlnymi,
a $q:V\to V/W$ jest homomorfizmem ilorazowym, to
obrazy $q(S_i)$ są zerowe albo nieprzywiedlne.
Z nich wybieramy maksymalną rodzinę, której suma jest
prosta. Jeżeli obraz $q(S_i)$ nie został wybrany i nie
można go do tej rodziny dołączyć, przecina on niezerowo
sumę wybranych obrazów. Nieprzywiedlność $q(S_i)$
sprawia, że cały ten obraz należy już do owej sumy.
Maksymalna suma zawiera więc wszystkie obrazy $q(S_i)$,
a zatem jest równa $V/W$, gdyż te obrazy generują iloraz.
Każdy niezerowy $q(S_i)$ jest obrazem
$S_i$ przez izomorfizm (jądro ograniczenia jest
podreprezentacją nieprzywiedlnej reprezentacji $S_i$).
Wybór odpowiednich $S_i$ daje więc podprzestrzeń
$U\subseteq V$ taką, że $q|_U:U\to V/W$ jest
izomorfizmem reprezentacji. Wtedy $V=W\oplus U$.


Dla implikacji $3\Rightarrow1$ wygodniejsza jest
kontrapozycja. Załóżmy $|G|\cdot1_K=0$.
Rozważmy reprezentację regularną na $V=K[G]$
z działaniem $\rho(h)([g])=[hg]$.
Wektor $u=\sum_{g\in G}[g]$ jest niezerowy,
gdyż w bazie $([g])_{g\in G}$ ma wszystkie współczynniki
równe $1_K$. Dla każdego $h\in G$ zachodzi
$\rho(h)u=u$, bo mnożenie z lewej permutuje elementy $G$.
Zatem $W=Ku$ jest jednowymiarową podreprezentacją.

Gdyby $W$ miała niezmiennicze dopełnienie $U$,
rzut $P:V=W\oplus U\to W$ wzdłuż $U$ byłby
$G$-równoważny: dla $w\in W$ i $v\in U$
mamy $P(\rho(h)(w+v))=\rho(h)w=\rho(h)P(w+v)$.
Ponieważ $[g]=\rho(g)[e]$ i działanie $G$ na $W$
jest trywialne, istnieje $a\in K$ takie, że
$P([g])=\rho(g)P([e])=au$ dla wszystkich $g\in G$.
Wobec tego
$P(u)=\sum_gP([g])=(|G|\cdot1_K)au=0$.
Ale $P|_W=\operatorname{id}_W$, więc $P(u)=u\ne0$:
sprzeczność.

Gdyby zachodził warunek 3, to reprezentacja regularna
byłaby sumą nieprzywiedlnych reprezentacji. W takim
rozkładzie każda podreprezentacja ma niezmiennicze
dopełnienie, co wynika ze struktury modułu półprostego.
Ponieważ dla konkretnej podreprezentacji $W$ wykazaliśmy
brak dopełnienia, warunek 3 również nie może zachodzić,
na mocy wykazanej już implikacji $3\Rightarrow2$.
Stąd $3\Rightarrow1$, co kończy dowód.
\end{proof}

\section{Dodatkowa kontrola założeń i przykład rozkładu}
\begin{uwaga}[Znaczenie charakterystyki]\label{aa:rem:maschke-char}
Zapis $\operatorname{char}K\nmid |G|$ w Twierdzeniu~\ref{aa:tw:maschke}
oznacza, że $|G|\cdot1_K\ne0$. Dla charakterystyki zero
warunek jest spełniony automatycznie; dla charakterystyki
pierwszej $p$ oznacza $p\nmid|G|$. Bez tego warunku
nie wolno dzielić przez $|G|$ w definicji operatora $P$.
Przykład~\ref{aa:prz:maschke-counter} pokazuje, że problem
nie jest wyłącznie techniczny: teza o dopełnieniu może zawieść.
\end{uwaga}
\begin{przyklad}[Rozkład reprezentacji permutacyjnej $C_2$]
\label{aa:prz:maschke-explicit}
Niech $G=C_2=\{e,s\}$, $K=\mathbb Q$ i niech
$\rho(s)(x,y)=(y,x)$ na $V=\mathbb Q^2$.
Wektory $u=(1,1)$ i $v=(1,-1)$ są liniowo niezależne,
a $\rho(s)u=u$, $\rho(s)v=-v$. Dlatego
$V=\mathbb Qu\oplus\mathbb Qv$ jest sumą dwóch
nieprzywiedlnych reprezentacji jednowymiarowych.
Dla dowolnego $(x,y)$ współczynniki tego rozkładu wynoszą
$(x+y)/2$ przy $u$ i $(x-y)/2$ przy $v$.
W charakterystyce $2$ te dwa wektory byłyby równe,
co unaocznia rolę odwracalności liczby $2$.
Możemy w tym przykładzie także \emph{obliczyć samo
uśrednianie} z dowodu Twierdzenia~\ref{aa:tw:maschke}.
Dla $W=\mathbb Q(1,1)$ wybieramy rzut
$p(x,y)=(x,x)$. Nie jest on przemienny z zamianą
współrzędnych: $\rho(s)p\rho(s)^{-1}(x,y)=(y,y)$.
Uśredniony rzut ma więc postać
\[
P(x,y)=\tfrac12\big((x,x)+(y,y)\big)
=\left(\tfrac{x+y}{2},\tfrac{x+y}{2}\right).
\]
Jego obrazem jest $W$, a jądrem jest
$\mathbb Q(1,-1)$. Otrzymaliśmy dokładnie dopełnienie
niezmiennicze, o którego istnieniu orzeka twierdzenie.
\end{przyklad}

\section{Średnia niezmiennicza i jej dokładny obraz}
\begin{stwierdzenie}[Operator uśredniania wektorów]\label{aa:prop:reynolds-vectors}
Niech $G$ będzie grupą skończoną, $K$ ciałem takim, że
$|G|\cdot1_K\ne0$, oraz $\rho:G\to\operatorname{GL}_K(V)$
reprezentacją na przestrzeni wektorowej $V$ nad $K$ (niekoniecznie
skończenie wymiarowej). Niech
$V^G=\{v\in V:\rho(g)v=v\ \text{dla każdego }g\in G\}$.
Wówczas operator
\[
R_G(v)=\frac1{|G|}\sum_{g\in G}\rho(g)v
\]
jest $K$-liniowym rzutem na $V^G$. W szczególności
$V=V^G\oplus\ker R_G$, a oba składniki są $G$-niezmiennicze.
\end{stwierdzenie}
\begin{proof}
Zbiór $V^G$ jest podprzestrzenią: suma dwóch wektorów
stałych oraz iloczyn wektora stałego przez skalar są stałe,
gdyż każdy $\rho(g)$ jest liniowy. Suma definiująca $R_G$
jest skończona, więc określa operator liniowy. Dla $h\in G$
i $v\in V$ mamy
\[
\rho(h)R_G(v)=\frac1{|G|}\sum_{g\in G}\rho(hg)v
=\frac1{|G|}\sum_{k\in G}\rho(k)v=R_G(v),
\]
gdyż mnożenie z lewej przez $h$ jest bijekcją $G$.
Wynika stąd $\operatorname{im}R_G\subseteq V^G$.
Jeżeli $v\in V^G$, wszystkie składniki sumy są równe $v$,
a więc $R_G(v)=v$. Zatem $\operatorname{im}R_G=V^G$
i $R_G^2=R_G$. Każdy $v\in V$ rozkłada się na
$R_G(v)+(v-R_G(v))$, przy czym drugi składnik należy
do jądra, ponieważ $R_G(v-R_G(v))=0$.
Przecięcie obrazu i jądra jest zerowe, ponieważ $R_G$
działa na swoim obrazie identycznościowo.
Wreszcie $R_G\rho(h)=R_G$ (podstawienie $k=gh$),
więc jądro jest niezmiennicze; obraz składa się z wektorów
stałych, a zatem również jest niezmienniczy.
\end{proof}

\begin{wniosek}[Dokładność brania wektorów niezmienniczych]
\label{aa:cor:invariants-exact}
Niech $G$ będzie grupą skończoną, $K$ ciałem takim, że
$|G|\cdot1_K\ne0$, a
\[
0\longrightarrow U\xrightarrow{\,i\,}V
 \xrightarrow{\,q\,}W\longrightarrow0
\]
krótkim ciągiem dokładnym reprezentacji $G$ nad $K$:
$i$ jest różnowartościowym operatorem przeplatającym,
$q$ jest surjektywnym operatorem przeplatającym oraz
$\operatorname{im}i=\ker q$. Wtedy ciąg
\[
0\longrightarrow U^G\xrightarrow{\,i|_{U^G}\,}V^G
 \xrightarrow{\,q|_{V^G}\,}W^G\longrightarrow0
\]
jest dokładny. Nie zakładamy skończonego wymiaru reprezentacji.
\end{wniosek}
\begin{proof}
Ponieważ $i$ i $q$ przeplatają działania grupy, przekształcają
wektory niezmiennicze w wektory niezmiennicze, toteż ograniczenia
w tezie są dobrze określone. Pierwsze jest różnowartościowe,
a dla $v\in V^G$ warunek $q(v)=0$ jest równoważny
$v\in\ker q=\operatorname{im}i$. Wówczas $v=i(u)$ dla
pewnego $u\in U$. Dla dowolnego $g\in G$ mamy
$i(\rho_U(g)u)=\rho_V(g)i(u)=v=i(u)$;
różnowartościowość $i$ daje $\rho_U(g)u=u$.
Stąd $\ker(q|_{V^G})=i(U^G)$.

Pozostaje surjektywność ostatniej strzałki. Niech $w\in W^G$.
Surjektywność $q$ pozwala wybrać $v\in V$ z $q(v)=w$.
Uśredniając według Stwierdzenia~\ref{aa:prop:reynolds-vectors},
otrzymujemy $R_G^V(v)\in V^G$. Ponieważ $q$ przeplata
działania, suma skończona spełnia
\[
q(R_G^V(v))=\frac1{|G|}\sum_{g\in G}q(\rho_V(g)v)
=\frac1{|G|}\sum_{g\in G}\rho_W(g)q(v)
=\frac1{|G|}\sum_{g\in G}w=w.
\]
Zatem każdy $w\in W^G$ ma niezmienniczy przeciwobraz,
co kończy dowód.
\end{proof}
\begin{przyklad}[Dlaczego branie niezmienników zawodzi w złej charakterystyce]
\label{aa:prz:invariants-not-exact-modular}
Niech $K=\mathbb F_2$ i $G=C_2=\{e,s\}$.
Rozważmy reprezentację regularną $K[G]$ z
Przykładu~\ref{aa:prz:group-algebra-modular} oraz homomorfizm
$G$-modułów $\varepsilon:K[G]\to K$ określony wzorem
$\varepsilon(a[e]+b[s])=a+b$, gdzie $G$ działa na $K$
trywialnie. Mapa $\varepsilon$ jest surjektywna, gdyż
$\varepsilon([e])=1$. Wektor $a[e]+b[s]$ jest niezmienniczy
wtedy i tylko wtedy, gdy $a=b$, zatem
$K[G]^G=K([e]+[s])$. Jednak
$\varepsilon([e]+[s])=1+1=0$ w $\mathbb F_2$, więc
$\varepsilon:K[G]^G\to K^G=K$ jest mapą zerową,
a nie surjekcją. Mimo dokładności ciągu
$0\to\ker\varepsilon\to K[G]\xrightarrow{\varepsilon}K\to0$
ciąg jego niezmienników nie jest dokładny na końcu.
\end{przyklad}

\begin{przyklad}[Średnia współrzędnych w reprezentacji permutacyjnej]
\label{aa:prz:reynolds-permutation}
Niech $G=S_3$ działa na $V=\mathbb Q^3$ przez permutowanie
współrzędnych. Dla $v=(a,b,c)$ każda z liczb $a,b,c$
występuje dokładnie dwa razy na każdej pozycji w sześciu
wektorach $\rho(g)v$. Dlatego
\[
R_G(a,b,c)=\left(\frac{a+b+c}{3},\frac{a+b+c}{3},
\frac{a+b+c}{3}\right).
\]
Obrazem jest prosta $\mathbb Q(1,1,1)$, a jądrem płaszczyzna
$a+b+c=0$. Są to dokładnie składniki rozkładu z
Przykładu~\ref{aa:prz:permutation-rep}. Warunek odwracalności
$|G|$ w ciele nie może być pominięty w ogólnym wzorze;
przykład nad $\mathbb F_2$ z
\ref{aa:prz:maschke-counter} pokazuje, że nawet istnienie
niezmienniczego dopełnienia może wtedy zawieść.
\end{przyklad}
